\subsection{A completely specified collar counterexample}

\begin{proposition}[Proved: fixed frequency and fixed collar width]\label{A2-collar-explicit}
Let $M$ be the hyperbolic pant with all three geodesic boundary lengths equal to $2\pi$. Label the boundaries $B_1,B_2,B_3$, and take an arclength coordinate $s\in\R/(2\pi\Z)$ on $B_1$, with origin at its shortest perpendicular to $B_2$. Set
\[
 n=10^4,\qquad \varepsilon=10^{-14},\qquad
 a_1(s)=1+\tfrac12\cos(ns),\qquad a_2=a_3=1.
\]
In inward Fermi coordinates remove the three collars $0\le r<\varepsilon a_i(s)$. They are disjoint, everywhere inward, and have analytic boundaries. On the retained pant let $g'=e^{2u}g$ be its conformally compatible hyperbolic metric with geodesic boundary. At the point of its first boundary with $s=\pi/n$ one has
\[
 u\le-\frac{9991}{20}\,10^{-14}<0.
\]
The same function $u$ is positive somewhere. Thus the general-shape counterexample of Proposition~\ref{A2-collar-shape} can be given with these fixed parameters, without an unspecified perturbation threshold.
\end{proposition}

\begin{proof}
Put $y=\arctan(\sinh r)$, so
\[
 g=\sec^2y\,(ds^2+dy^2),\qquad
 \Delta_g=\cos^2y\,(\partial_s^2+\partial_y^2),\qquad
 L=\Delta_g-2.
\]
The collar theorem, applied to the double of $M$, gives disjoint collars of Fermi width $\operatorname{arsinh}(1/\sinh\pi)$; see \cite[Theorem~2.18, pp.~22--23]{McMullen1994A567}. In particular the three closed collars $0\le y\le1/50$ are embedded and disjoint. Indeed $\pi<4$, $e<3$, and
\[
 \tan(1/50)\le\frac{1/50}{1-(1/50)^2/2}
 =\frac{100}{4999}<\frac2{81}<\frac1{\sinh\pi}.
\]
The selected new boundaries lie in these collars, and even in $y<1/100$, since $0<\varepsilon a_i\le3\varepsilon/2<1/100$ and $y\le r$.

Here are explicit comparison functions. Write $b=1/100$, and define the $C^2$ cutoff
\[
 \eta(y)=\begin{cases}
 1,&0\le y\le b,\\
 1-10t^3+15t^4-6t^5,&b\le y\le2b,\quad t=(y-b)/b,\\
 0,&y\ge2b.
 \end{cases}
\]
Its first two derivatives vanish at the endpoints of the transition interval. On that interval,
\[
 |\eta'|\le\frac{15}{8b}<200,\qquad
 |\eta''|\le\frac{10}{\sqrt3\,b^2}<60000.
\]
These bounds follow by differentiating the displayed polynomial: its derivatives are $-30t^2(1-t)^2$ and $-60t(1-t)(1-2t)$. Their exact maximum absolute values on $[0,1]$ are $15/8$ and $10/\sqrt3$. Also, throughout $0\le y\le2b$,
\[
 \tan y<\frac1{40},\qquad \sec^2y<2.
\]
For the second inequality use $\cos y\ge1-y^2/2\ge4999/5000$.

Extend each $\eta(y_i)\tan y_i$ by zero outside its collar and set
\[
 H=2000-\sum_{i=1}^3\eta(y_i)\tan y_i.
\]
At most one summand is nonzero at a point. Since $L(\tan y)=0$, differentiation gives
\[
 LH=-4000-\cos^2y\bigl(\eta''\tan y+2\eta'\sec^2y\bigr)
 \le-4000+60000/40+800=-1700
\]
on a collar, and $LH=-4000$ elsewhere.

The separated function
\[
 f_n(y)=e^{-ny}(n-\tan y)
\]
satisfies the exact identity $f_n''-(n^2+2\sec^2y)f_n=0$ and $-f_n'(0)=n^2+1$. Define
\[
 w=\eta(y_1)f_n(y_1)\cos(ns)
\]
on the first collar and extend it by zero. Its residual is supported in $b\le y_1\le2b$ and equals
\[
 Lw=\cos^2y_1\bigl(\eta''f_n+2\eta'f_n'\bigr)\cos(ns).
\]
Consequently, using $nb=100$ and $e>2$,
\begin{align*}
 \|Lw\|_\infty
 &\le e^{-100}\left[60000(n+1/40)
       +400(n^2+n/40+2)\right]\\
 &<\frac{40600102300}{2^{100}}<1.
\end{align*}
All bounds so far are explicit rational or integer inequalities apart from the elementary displayed trigonometric estimates.

Set $\Psi=2H+w/2$ and $U=\varepsilon\Psi$. Then $L\Psi\le-3400+1/2<0$. Since $e^{2x}-1\ge2x$ for every real $x$, this proves
\[
 \Delta_g U\le e^{2U}-1.
\]
It remains to check the outward normal derivative on each new boundary; an interior supersolution alone would not suffice.

Write a boundary graph as $r=f(s)=\varepsilon a(s)$, and put $p=f'(s)$, $c=\cosh f$, $q=p/c$. Its outward normal and signed geodesic curvature are exactly
\[
 N=\frac{-c\,\partial_r+(p/c)\partial_s}{\sqrt{c^2+p^2}},\qquad
 k_g=\frac{\tanh f(1+2q^2)-f''\operatorname{sech}^2f}
                {(1+q^2)^{3/2}}.
\]
For example, the curvature formula follows by inserting the graph velocity $f'\partial_r+\partial_s$ into $g(\nabla_TT,N)$, using $\Gamma^r_{ss}=-\cosh r\sinh r$ and $\Gamma^s_{rs}=\tanh r$.

We give the error bounds used in the comparison. On the region $0\le r\le1/100$, the cutoff is one. On the first collar there,
\[
 \Psi=4000-2\sinh r+\tfrac12F_n(r)\cos(ns),\qquad
 F_n(r)=e^{-n\arctan(\sinh r)}(n-\sinh r).
\]
On the other two collars the oscillating term is absent. Because $\sinh r<1$, $\cosh r<2$, $|y'|\le1$, and $|y''|\le1$, direct differentiation gives
\[
 |F_n|\le2n,\quad |F_n'|\le2n^2,\quad
 |F_n''|\le n^3+2n^2+5n+1\le4n^3.
\]
It follows that
\[
 |\Psi_r|\le2n^2,\qquad |\Psi_s|\le n^2,
 \qquad |\Psi_{rr}|\le3n^3.
\]
At $r=0$ the outward derivative is
\[
 -\Psi_r(0,s)=a(s)-a''(s)+1.
\]
The inequalities $f\le3\varepsilon/2$, $|p|\le\varepsilon n/2$, and
$1-(1+q^2)^{-1/2}\le q^2/2$ therefore imply
\begin{align*}
 |\partial_N\Psi-(a-a''+1)|
 &\le 3n^3f+n^2p^2+n^2|p|\\
 &\le5\varepsilon n^3+\tfrac14\varepsilon^2n^4
 <10\varepsilon n^3.
\end{align*}
The same estimates apply on the other two boundary components.

For the curvature, use $|\tanh f-f|\le f^3/3$,
$1-\operatorname{sech}^2f\le f^2$, and
$1-(1+x)^{-3/2}\le3x/2$ for $x\ge0$. On $0\le x\le1$ the derivative of $(1+2x)(1+x)^{-3/2}$ has absolute value at most $1/2$, so its distance from one is at most $2x$. Since $q^2\le p^2<1$, the exact curvature formula yields
\begin{align*}
 |k_g-\varepsilon(a-a'')|
 &\le f^3/3+2fq^2+|f''|(f^2+3q^2/2)\\
 &\le\varepsilon^3\left(\tfrac98+\tfrac{15}8n^2+\tfrac3{16}n^4\right)
 \le\varepsilon^3n^4.
\end{align*}
Thus, for the specified $n,\varepsilon$,
\[
 \partial_NU-k_g
 \ge\varepsilon\bigl(1-10\varepsilon n^3-\varepsilon^2n^4\bigr)>0,
 \qquad 10\varepsilon n^3+\varepsilon^2n^4=\frac1{10}+10^{-12}<1.
\]

The function $u$ satisfies $\Delta_gu=e^{2u}-1$ and $\partial_Nu=k_g$. Apply the maximum principle and boundary-point lemma to $U-u$: it satisfies $(\Delta_g-c_0)(U-u)\le0$ with $c_0>0$ and outward derivative nonnegative. A negative minimum is impossible, whether interior or on the boundary. Hence $u\le U$ on the retained pant.

At the specified boundary point, $f=\varepsilon/2$ and $\cos(ns)=-1$. Here $ny\le n\varepsilon/2<1/10$, so $e^{-ny}\ge1-ny\ge9/10$, while $\sinh f<1$. Therefore
\[
 \Psi\le4000-\tfrac9{20}(n-1)=-\frac{9991}{20}.
\]
This proves the stated strict negative bound. The new boundaries are analytic because their defining functions $a_i$ are analytic; the merely $C^2$ cutoff occurs only in an auxiliary comparison function. Their positive inward widths and the explicit disjoint collars verify every geometric hypothesis. Finally the retained surface has old area less than $2\pi$, while its new geodesic-boundary hyperbolic area is $2\pi$ by Gauss--Bonnet. Hence $u>0$ somewhere. The exact symbolic identities and rational inequalities used above are checked in \path{code/part_a567_collar_explicit.py}; no numerical approximation to $u$ is used.
\end{proof}
